\section{Evaluation} \label{sec:eval}

We discusses the experience in using \sysname{} and
its performance in checking compatibilities.

\newcommand{\tableT}{\textcolor{blue}{T}}
\newcommand{\tableF}{\textcolor{red}{F}}
\newcommand{\tableNA}{-}
\begin{table}[H]
  \centering
  \begin{tabular}{|l|c|c|c|c|c|c|c|c|}
    \hline
    {} &
    EC &
    MR &
    MW &
    RYW &
    WFR &
    PRAM &
    CC &
    LIN \\
    \hline
    EC   & \tableNA & \textcolor{red}{24} & \textcolor{red}{22} &
    \textcolor{red}{26} & \textcolor{red}{25} & \textcolor{red}{24} &
    \textcolor{red}{35} & \textcolor{red}{32} \\
    MR   & \textcolor{blue}{23} & \tableNA & \textcolor{red}{12} &
    \textcolor{red}{16} & \textcolor{red}{75} & \textcolor{red}{16} &
    \textcolor{red}{26} & \textcolor{red}{217} \\
    MW   & \textcolor{blue}{27} & \textcolor{red}{16} & \tableNA &
    \textcolor{red}{18} & \textcolor{red}{16} & \textcolor{red}{16} &
    \textcolor{red}{23} & \textcolor{red}{22} \\
    RYW  & \textcolor{blue}{193} & \textcolor{red}{106} &
    \textcolor{red}{45} & \tableNA & \textcolor{red}{59} &
    \textcolor{red}{20} & \textcolor{red}{29} & \textcolor{red}{26} \\
    WFR  & \textcolor{blue}{24} & \textcolor{red}{125} &
    \textcolor{red}{12} & \textcolor{red}{17} & \tableNA &
    \textcolor{red}{18} & \textcolor{red}{44} & \textcolor{red}{54} \\
    PRAM & \textcolor{blue}{24} & \textcolor{blue}{14} &
    \textcolor{blue}{10} & \textcolor{blue}{14} &
    \textcolor{red}{178} & \tableNA & \textcolor{red}{27} &
    \textcolor{red}{28} \\
    CC   & \textcolor{blue}{33} & \textcolor{blue}{20} &
    \textcolor{blue}{16} & \textcolor{blue}{19} &
    \textcolor{blue}{19} & \textcolor{blue}{20} & \tableNA &
    \textcolor{red}{282} \\
    LIN  & \textcolor{blue}{26} & \textcolor{blue}{21} &
    \textcolor{blue}{17} & \textcolor{blue}{22} &
    \textcolor{blue}{22} & \textcolor{blue}{22} &
    \textcolor{blue}{44} & \tableNA \\
    \hline
  \end{tabular}
  \caption{Pairwise check performance. Semantics in the first column
    are checked if they satisfy semantics in other columns. Each cell
    shows the execution time for completing the check in milliseconds. The
  blue text indicates the result of the check is true and red means false. }
  \label{tab:pairwise-time}
\end{table}
\vspace{-0.3in}

\para{Experience}
We implemented \sysname{} in Python on Z3 version 4.15.1. It took
approximately 3.7K lines of code to implement \sysname{}. 
%
%% \textbf{Usability}
%% During review, on several instances we were asked on the usability of \sysname.
%% We summarized them in these 3 main categories: how can the framework be used,
%% the expected user input and effort,
%% and the scalability of our methodology in large-scale applications.
% 1
To use \sysname{}, users supply \sysgraph whose nodes and/or edges carry
built-in or custom \sysmods that extend the base model and relation interfaces.
When multiple \sysmods are attached, \sysname checks for compatibility of given
input or exhaustively explores valid combinations.

%The \texttt{check} function optionally emits proof goals via the \texttt{witness} argument (see \texttt{tests/test\_witness.py} in the accompanying artifact), which can be discharged to external SMT solvers.

% 2
On effort and learning curve, once semantic requirements are understood, writing \sysgraph is direct.
Our case studies only require 183, 80, and 252 lines of code for encoding the moving streaming service, online shopping mall, and XC storage, respectively, thanks to existing facilities such as the \sysmods{} in the \sysname{} framework. 
These relatively small efforts for constraint encoding and automated checks justifies our framework's intention for agile compatibility checks compared to typical end-to-end formal verification which could have 20:1 proof-to-code/specification ratio~\cite{certikos-osdi16, klein2009sel4, wormspace, ironfleet} which could follow with confidence after our check passes. 
%Our effort constraints remain stable unless interfaces or consistency requirements change.

% 3
The scalability of our tool is mainly expressed via modular composition,
\sysgraph nodes and edges compose modularly: large systems can be decomposed into sub-\sysgraphs,
verified independently, and re-aggregated into monolithic node that preserves the verified properties of its components.
This pattern (\eg Case Study 2 in Section \ref{sec:cs}) also enables parallel SMT solving on sub-components and 
reuse of solved results for large service meshes.
%(see \texttt{tests/test\_aggregate.py} and Case Study 2 in Section \ref{sec:cs}) also enables parallel SMT solving on sub-components and reuse of solved results for large service meshes.

% end
In general, our framework offers early compatibility detection,
architecture validation before implementation, and pre-verification screening to de-risk downstream formal proofs.
It targets teams building correctness-critical distributed systems who want a high-level, low-maintenance checker rather than line-by-line deductive proof development (\eg using tools like Verus\cite{verus}, Dafny\cite{dafny}, Rocq\cite{coq}, and Lean\cite{lean}).


\para{Performance}
SMT solvers often take an extended period to complete their tasks.
Thus, we measure the
duration of the compatibility check in \sysname{}.
All experiments were conducted on an Intel i7-1360P processor with
64GB RAM running NixOS 25.11 (Linux 6.12.30).

% 10 seconds
% timeout limit. However, during our evaluation, we discovered a
% termination bug in Z3 4.15.1
% where the solver would not respect timeout parameters, continuing to
% run well beyond the specified limit (e.g., still running after 2
% minutes despite a 10-second timeout). To address this, we implemented
% a watchdog mechanism using strace to monitors write systemcalls to
% stdout, sending SIGINT if the 10-second timeout is exceeded,
% following the philosophy established in recent
% verification work Verus that ``it's better to improve solver
% performance than rely on a slow proof'' \cite{verus}.
% https://github.com/verus-lang/verus/blob/0b6f3cbb8fceff6681b4988c3cac37a8b801f8e6/source/docs/guide/src/smt_failures.md

\textbf{Pairwise Compatibility Check.}

First, we evaluated the compatibility between pairs of base
consistency semantics,
EC, MR, RYW, MW, WFR, PRAM, CC, and LIN, in \sysname{}.
Table \ref{tab:pairwise-time} shows the execution times for these
pairwise compatibility checks and the results of the check encoded in
colors (blue: compatible, red: non-compatible). On average, it took
only 31 ms per check and the results of the checks are correct with
respect to the theory.

\textbf{Compatible \sysmod{} Search.} Next, we evaluate \sysname{} using
three case studies in \cref{sec:cs}.
This evaluation focuses on two key metrics: the framework's ability to
find valid compositions with a minimal semantic score and its
performance characteristics under different system complexities.

For each case study, we gave blank \sysmods{} to each application
so that \sysname{} search and output compatible \sysmod{} combinations.
\sysname{} uses the search strategy in \cref{subsec:search}.
%we enumerated all possible combinations of
%consistency models for nodes with non-null provs fields, while
%keeping nodes with predefined needs fields unchanged. We assigned
%scores to individual consistency models and their combinations based
%on their strength as shown in Figure \ref{fig:semantics-tree}.

Movie streaming service requires finding \sysmods{} corresponding to
five DB nodes that satisfy application node and edge constraints.
Given that \sysname{} contains 17 \sysmods{} as illustrated in
\cref{fig:semantics-tree} the full search space contains
$17^5$ combinations. On average checking each combination
takes 16.5 ms and can find a compatible combination with the lowest
score in 20.43 seconds with a timeout of 10 seconds for each check.
Checking the entire search space is feasible but is expected to take
64.82 hours (we have completed 800,507 checks or 56.28\% of search space
in 36.55 hours).

The online shopping mall example is simpler than the movie streaming service
and involves only three DB nodes to search.
We evaluated all $17^3$ possible combinations of consistency semantics across
these nodes. The evaluation completed in 251.77 seconds with an average
duration of 5.12 ms per combination with 0 timeouts.

Finally, four DB node combinations (i.e., $17^4$ cases) are searched for
the cross causal storage. The total search completes in 3.98 hours
with an average duration of 17.17 ms per combination.
Our tool successfully identified the combination in our original design
in about 261.09 seconds. The timeout rate is only about 0.15\% of cases.

Overall, depending on the complexity of the constraints the
compatibility check time
varies, but individual checks are done in less than 20ms on average. Even though
searching through compatible answers in the entire search space could
take extended
time, finding the most desired, minimum score answer for each case
study takes only less than 5 minutes.
